%% kapa: exam=1 round=28 year=2017 period=1 form=A subject=real_estate subject_ko=부동산학원론 seq=20 number=100
%% source: Q-Net 감정평가사 기출문제 (cst00309) artlSeq=5199603 file=extracted/1008/5199603/2017년 제28회 감정평가사 1차시험 문제지/1교시 A형.pdf
%% conversion: kordoc 4.24.51 -> tools/gen1.py (auto); review=proofread-sonnet; flags=none
\begin{question}{100}
부동산 권리분석에 관한 내용으로 옳지 않은 것은?
\begin{choices}
  \item 부동산의 상태 또는 사실관계, 등기능력 없는 권리 및 등기를 요하지 않는 권리관계 등 자세한 내용에 이르기까지 분석의 대상으로 하는 것이 협의의 권리분석이다.
  \item 매수인이 대상부동산을 매수하기 전에 소유권이전을 저해하는 조세체납, 계약상하자 등을 확인하기 위해 공부 등을 조사하는 일도 포함된다.
  \item 부동산 권리관계를 실질적으로 조사, 확인, 판단하여 일련의 부동산활동을 안전하게 하려는 것이다.
  \item 대상부동산의 권리에 하자가 없는지 여부를 판단하는 것을 권리분석이라 한다.
  \item 권리분석 보고서에는 대상부동산 및 의뢰인, 권리분석의 목적, 판단결과의 표시 및 이유, 권리분석의 방법 및 성격, 수집한 자료 목록, 면책사항 등이 포함된다.
\end{choices}
\end{question}
\answer{1}
